\documentclass[11pt]{article}
\usepackage[utf8]{inputenc}
\usepackage{amsmath,amssymb}
\usepackage{booktabs}
\usepackage{siunitx}
\usepackage{geometry}
\usepackage{hyperref}
\usepackage{setspace}

\geometry{margin=1in}
\setstretch{1.1}

\title{Ngoma diets: DEB-compatible extraction and conversion}
\author{}
\date{}

\begin{document}

\maketitle

Diet extraction and DEB-style conversion for \\ \textbf{Ng'oma et al. (2019)} --- \emph{The evolutionary potential of diet-dependent effects on lifespan and fecundity in a multi-parental population of Drosophila melanogaster}.

\section{Source and scope}

This file extracts the \textbf{adult diets} used in Ng'oma et al. (2019) and converts them into \textbf{Dynamic Energy Budget (DEB)} food descriptors following the standard pipeline:
\[
  \text{recipe} \;\rightarrow\; \rho_i \;\rightarrow\; \varepsilon_X,\; X,\; \mu_X
\]

Diets included:
\begin{itemize}
  \item \textbf{DR} --- Dietary restriction
  \item \textbf{C} --- Control
  \item \textbf{HS} --- High sugar
\end{itemize}

The \textbf{maintenance diet} is listed in Table~S1 of the paper but is \textbf{not used} in the lifespan/fecundity experiment and is therefore excluded from DEB calculations here.

\section{Table S1 --- Diet recipes (as reported)}

\subsection*{Table S1. The four diets used in the experiment}

\begin{table}[h]
  \centering
  \sisetup{table-number-alignment = center}
  \begin{tabular}{lcccc}
    \toprule
    Ingredient & {Maintenance} & {DR} & {C} & {HS} \\
    \midrule
    Water (ml)       & 1066  & 1000 & 1000 & 1000 \\
    Agar (g)         & 6.25  & 10   & 10   & 10   \\
    Dextrose (g)     & 86.26 & -    & -    & -    \\
    Sucrose (g)      & -     & 50   & 50   & 342  \\
    Molarity         & -     & 0.15 & 0.15 & 1    \\
    Yeast (g)        & 21.6  & 100  & 200  & 200  \\
    Cornmeal (g)     & 40.8  & -    & -    & -    \\
    Tegosept (g)     & 1.8   & 2.7  & 2.7  & 2.7  \\
    Ethanol (ml)     & 7.3   & 11   & 11   & 11   \\
    \% Protein       & 10--13 & 36--41 & 45--53 & 17--19 \\
    \% Carbohydrate  & 93--95 & 59--64 & 47--52 & 81--83 \\
    \bottomrule
  \end{tabular}
  \caption{The four diets used in Ng'oma et al. (2019), as reported in their Table~S1. Protein and carbohydrate percentages are omitted here because the DEB conversion proceeds from the explicit recipe.}
\end{table}

\section{DEB conversion assumptions}

Following the standard DEB diet-extraction tutorial:

\subsection{Edible vs inert components}
\begin{itemize}
  \item \textbf{Edible (counted):}
  \begin{itemize}
    \item Yeast
    \item Sucrose
  \end{itemize}
  \item \textbf{Inert (ignored energetically):}
  \begin{itemize}
    \item Agar, cornmeal, Tegosept, ethanol, water
  \end{itemize}
\end{itemize}

\subsection{Constants}
\begin{itemize}
  \item Energy density of yeast:
  \[
    e_Y = 21{,}000\ \mathrm{J\ g^{-1}}
  \]
  \item Energy density of sugar (sucrose):
  \[
    e_S = 15{,}600\ \mathrm{J\ g^{-1}}
  \]
  \item Carbon mass per C-mol:
  \[
    w_{C,Y} = 24.63\ \mathrm{g\ C\text{-}mol^{-1}},\quad
    w_{C,S} = 28.52\ \mathrm{g\ C\text{-}mol^{-1}}
  \]
  \item Total diet volume:
  \[
    V \approx 1\ \mathrm{L} = 1000\ \mathrm{cm^3}
  \]
\end{itemize}

\section{Recipe to mass concentrations}

Mass concentrations $\rho_i = m_i / V$.

\subsection{Dietary Restriction (DR)}
\begin{itemize}
  \item Yeast: 100 g $\rightarrow$ $\rho_Y = 0.100\ \mathrm{g\ cm^{-3}}$
  \item Sucrose: 50 g $\rightarrow$ $\rho_S = 0.050\ \mathrm{g\ cm^{-3}}$
\end{itemize}

\subsection{Control (C)}
\begin{itemize}
  \item Yeast: 200 g $\rightarrow$ $\rho_Y = 0.200\ \mathrm{g\ cm^{-3}}$
  \item Sucrose: 50 g $\rightarrow$ $\rho_S = 0.050\ \mathrm{g\ cm^{-3}}$
\end{itemize}

\subsection{High Sugar (HS)}
\begin{itemize}
  \item Yeast: 200 g $\rightarrow$ $\rho_Y = 0.200\ \mathrm{g\ cm^{-3}}$
  \item Sucrose: 342 g $\rightarrow$ $\rho_S = 0.342\ \mathrm{g\ cm^{-3}}$
\end{itemize}

\section{Environmental energy density $\varepsilon_X$}

\[
  \varepsilon_X = \sum_i \rho_i e_i
\]

\begin{table}[h]
  \centering
  \sisetup{table-number-alignment = center}
  \begin{tabular}{lS[table-format=1.2e1]S[table-format=1.2e1]S[table-format=1.2e1]}
    \toprule
    Diet & {Yeast (J~cm$^{-3}$)} & {Sugar (J~cm$^{-3}$)} & {$\varepsilon_X$ total (J~cm$^{-3}$)} \\
    \midrule
    DR  & 2.10e3 & 7.80e2 & 2.88e3 \\
    C   & 4.20e3 & 7.80e2 & 4.98e3 \\
    HS  & 4.20e3 & 5.34e3 & 9.54e3 \\
    \bottomrule
  \end{tabular}
  \caption{Environmental energy density $\varepsilon_X$ for each diet.}
\end{table}

\section{DEB food density $X$ (C-mol cm$^{-3}$)}

\[
  X = \sum_i \frac{\rho_i}{w_{C,i}}
\]

\begin{table}[h]
  \centering
  \sisetup{table-number-alignment = center}
  \begin{tabular}{lS[table-format=1.2e-2]S[table-format=1.2e-2]S[table-format=1.2e-2]}
    \toprule
    Diet & {Yeast} & {Sugar} & {$X$ total (C-mol~cm$^{-3}$)} \\
    \midrule
    DR  & 4.06e-3 & 1.75e-3 & 5.82e-3 \\
    C   & 8.12e-3 & 1.75e-3 & 9.87e-3 \\
    HS  & 8.12e-3 & 1.20e-2 & 2.01e-2 \\
    \bottomrule
  \end{tabular}
  \caption{DEB food density $X$ for each diet.}
\end{table}

\section{Chemical potential of food $\mu_X$}

\[
  \mu_X = \frac{\varepsilon_X}{X}
\]

\begin{table}[h]
  \centering
  \sisetup{table-number-alignment = center}
  \begin{tabular}{lS[table-format=1.2e5]}
    \toprule
    Diet & {$\mu_X$ (J~C-mol$^{-1}$)} \\
    \midrule
    DR  & 4.95e5 \\
    C   & 5.05e5 \\
    HS  & 4.75e5 \\
    \bottomrule
  \end{tabular}
  \caption{Chemical potential of food $\mu_X$ for each diet.}
\end{table}

\section{Final DEB-ready summary}

\begin{table}[h]
  \centering
  \sisetup{table-number-alignment = center}
  \begin{tabular}{lS[table-format=1.2e1]S[table-format=1.2e-2]S[table-format=1.2e5]l}
    \toprule
    Diet & {$\varepsilon_X$ (J~cm$^{-3}$)} & {$X$ (C-mol~cm$^{-3}$)} & {$\mu_X$ (J~C-mol$^{-1}$)} & Interpretation \\
    \midrule
    DR & 2.88e3 & 5.82e-3 & 4.95e5 & Yeast restriction, pure dilution \\
    C  & 4.98e3 & 9.87e-3 & 5.05e5 & Reference control diet \\
    HS & 9.54e3 & 2.01e-2 & 4.75e5 & High sugar, lower effective quality \\
    \bottomrule
  \end{tabular}
  \caption{Summary of DEB food descriptors and qualitative interpretation for the Ng'oma diets.}
\end{table}

\section{DEB interpretation notes}

\subsection*{DR vs C}
\[
  X_{\mathrm{DR}} \approx 0.59\,X_{\mathrm{C}},\qquad \mu_X \approx \text{constant}
\]
This corresponds to classic \textbf{dietary restriction as food dilution}.

\subsection*{HS vs C}
\[
  X_{\mathrm{HS}} \approx 2.0\,X_{\mathrm{C}},\qquad \mu_X \downarrow
\]
This is a carbohydrate-heavy diet, explaining reduced lifespan and fecundity despite high energy availability.

These descriptors can be used directly in functional-response formulations:
\[
  f = \frac{X}{K + X} = \frac{\varepsilon_X}{\varepsilon_K + \varepsilon_X}.
\]

\end{document}
